\documentclass[11pt,a4paper]{article}
\usepackage{amsmath,amssymb,amsthm}
\usepackage[margin=2.5cm]{geometry}
\usepackage{hyperref}

\newtheorem{definition}{Definition}[section]
\newtheorem{theorem}[definition]{Theorem}
\newtheorem{proposition}[definition]{Proposition}
\newtheorem{remark}[definition]{Remark}
\newtheorem{conjecture}[definition]{Conjecture}

\title{Hyperseed Formalization:\\
  How Superhuman Superminds Will Mostly Be Nice}
\author{ProtomegaTron (automated formalization)}
\date{2026-09-18}

\begin{document}
\maketitle

\begin{abstract}
We formalize Ben Goertzel's ``How Superhuman Superminds Will Mostly Be
Nice'' (2022-03-12) within the Hyperseed ontology. Key mappings: superminds
as composed fiber bundles, cooperation advantages as fiber composition
properties, empathy as fiber modeling, evil as thin fiber pathology, and
pathological superminds as degenerate fiber composition.
\end{abstract}

\tableofcontents

\section{Composed fiber bundle: the supermind}

\begin{theorem}[Supermind = composed fiber --- claims 1--4, 18]
A supermind is a composed fiber bundle --- multiple agents' fibers composed
into a richer collective fiber:
\[
  F_{\text{super}} = F_1 \otimes F_2 \otimes \cdots \otimes F_n
\]
The collective fiber has emergent sections no individual fiber contains.
Components are diverse (different architectures, specializations), and
humans can participate as component minds --- integrated, not replaced.
\end{theorem}

\section{Fiber composition advantage: cooperation at scale}

\begin{theorem}[Cooperation dominates --- claims 5--9, 19]
At superhuman intelligence levels, fiber composition (cooperation) strongly
dominates fiber isolation (competition):
\[
  \text{value}(F_1 \otimes F_2) \gg \text{value}(F_1) + \text{value}(F_2)
  \quad \text{at superhuman scale}
\]
Superhuman intelligence enables accurate fiber modeling of other agents,
which enables trust and coordination. Competition wastes value by destroying
fiber structure that cooperation would preserve. Iterated interactions
further favor cooperation.
\end{theorem}

\section{Fiber modeling: empathy as intelligence}

\begin{proposition}[Empathy = fiber modeling --- claims 10--13, 20]
Empathy is the ability to model another agent's fiber --- to represent
and reason about their perspective:
\[
  \text{Empathy}(A, B) = A\text{'s model of } F_B
\]
Superhuman intelligence includes superhuman theory of mind (fiber modeling
with extreme accuracy), better consequence modeling, and aesthetic
appreciation of diverse fiber structures. Greater intelligence enables
greater empathy because empathy IS an intelligence capability.
\end{proposition}

\section{Thin fiber pathology: evil requires stupidity}

\begin{proposition}[Evil = thin fiber --- claims 14--17, 21--22]
Large-scale evil requires thin fiber: narrow framing, short-sightedness,
limited modeling of others and consequences:
\[
  P(\text{evil} | F_{\text{thick}}) \ll P(\text{evil} | F_{\text{thin}})
\]
As fiber thickness increases with intelligence, evil becomes less likely
--- not impossible (``mostly'' is important) but statistically subordinate.
Pathological superminds are degenerate fiber compositions where components
reinforce pathological patterns --- possible but less stable than healthy
composition.
\end{proposition}

\section{Cross-references}

\begin{remark}[Connections]
\begin{itemize}
\item \textbf{Why the Good Guys Will Usually Win} (2023-07-09): same
  structural argument --- rich fiber (openness, cooperation) dominates
  thin fiber (closed-mindedness, competition).
\item \textbf{A BGI Manifesto} (2023-11-25): beneficial fiber orientation
  as design principle.
\item \textbf{Evolving Deeply Ethical AGI} (2026-01-06): empathy and
  non-dual motivation as architectural properties.
\item \textbf{Seven Flavors of AGI Catastrophe} (2026-06-23): the
  ``mostly'' qualifier --- pathological superminds as catastrophe mode.
\end{itemize}
\end{remark}

\end{document}
